\documentclass[runningheads]{llncs}
\usepackage[T1]{fontenc}
\usepackage[utf8]{inputenc}
\usepackage{graphicx}
\usepackage{booktabs}
\usepackage{float}
\usepackage{amsmath}
\usepackage{url}
\usepackage[hidelinks]{hyperref}
\renewcommand{\tablename}{Tabla}
\renewcommand{\figurename}{Figura}
\renewcommand{\andname}{y}
\emergencystretch=2.5em
\begin{document}

\title{Taller 01 --- Foundation Models: Comparación, Decodificación y Razonamiento}
\titlerunning{Taller 01 --- Foundation Models}
\author{Jaime Astudillo \and Roberth Chachalo}
\authorrunning{J. Astudillo y R. Chachalo}
\institute{MMIA 6013 --- IA Generativa y Agentes, Universidad San Francisco de Quito}
\maketitle

\section{Tarea elegida}
Elegimos diez problemas de aritmética con respuesta entera para verificar
automáticamente los resultados. Agregamos tres casos contaminados y comparamos cómo
cambian las respuestas al modificar el modelo, el prompt y el esfuerzo.

Los diez casos limpios y los tres contaminados, con sus respuestas esperadas, están en \texttt{datos/casos.json} del repositorio.

\section{Parte 0 --- El modelo base y la distribución}
Fila \texttt{base\_local}: \texttt{openai-community/gpt2}, 124\,M de parámetros, en CPU, sin API ni credenciales.

\subsection{La distribución del siguiente token}
En el dominio de números y cálculo usamos \texttt{As easy as one, two,} como frase
hecha relativamente predecible y \texttt{To calculate the total, we need to} como
contexto abierto. Los prefijos están en inglés, la lengua predominante del
modelo base.
La certeza se comprueba con la probabilidad máxima y la entropía; no se supone.

Se hace un forward por prefijo y se reutilizan sus logits para todas las
temperaturas. La softmax, la entropía y el núcleo se calculan sobre \textbf{todo el
vocabulario}, aunque la primera figura solo muestre los 15 primeros tokens.

Al bajar la temperatura, observamos más probabilidad concentrada en los primeros
tokens. Al subirla, aumentaron la entropía y las alternativas. Con temperatura
1, el núcleo del 90 \% incluyó cinco tokens en el prefijo seguro y 241 en el
abierto.

\begin{figure}[H]\centering\includegraphics[width=0.95\textwidth,height=0.42\textheight,keepaspectratio]{figuras/parte0a_temperatura.png}\caption{Distribución de los 15 candidatos más probables para $T\in\{0.1,0.7,1.0,1.5,2.0\}$ en los dos prefijos.}\label{fig:p0a}\end{figure}
\begin{table}[H]\centering\caption{Entropía y tamaño del núcleo del 90\,\% por prefijo y temperatura (GPT-2).}\label{tab:p0}
{\footnotesize\setlength{\tabcolsep}{3.5pt}
\begin{tabular}{lrrrr}
\hline
prefijo & T & entropía (bits) & núcleo 0.9 & p máx \\
\hline
seguro & 0.1 & 3.444e-05 & 1 & 1 \\
seguro & 0.7 & 0.8858 & 2 & 0.8365 \\
seguro & 1 & 2.323 & 5 & 0.6482 \\
seguro & 1.5 & 9.068 & 7659 & 0.2155 \\
seguro & 2 & 13.23 & 20702 & 0.04256 \\
abierto & 0.1 & 0.0684 & 1 & 0.9937 \\
abierto & 0.7 & 5.129 & 40 & 0.1755 \\
abierto & 1 & 7.1 & 241 & 0.08293 \\
abierto & 1.5 & 9.894 & 1735 & 0.02624 \\
abierto & 2 & 11.82 & 6464 & 0.009824 \\
\hline
\end{tabular}
}\end{table}


\subsection{Las tres palancas sobre esa misma distribución}
Usamos \texttt{model.generate(...)}, modo evaluación y semilla registrada por corrida.
Para aislar cada palanca desactivamos los otros filtros: \texttt{top\_k=0}, \texttt{top\_p=1},
sin penalización de repetición y con un solo beam. El aviso de Transformers
al enviar temperatura en modo greedy se conserva: ese parámetro se ignora.
Véase la documentación de generación.

Cada generación se añade a \texttt{datos/resultados.csv}, con \texttt{parte=0} y costo cero.
\texttt{datos/parte0\_generaciones.jsonl} conserva además el prompt, los IDs, la semilla,
la revisión del modelo, la fecha y los parámetros exactos. En esta parte
\texttt{acierto} queda vacío porque no se evalúa exactitud.

Arriba se conservan las dos salidas del interruptor,
las cinco corridas de \texttt{top\_k=1} con sus IDs y soportes, la figura de cortes
y la continuación de cien tokens sin editar. El conteo de cuatrigramas
repetidos complementa la observación del texto; no evalúa su corrección.

Salida con \texttt{do\_sample=False} (idéntica para $T=0.2$ y $T=1.5$):
\begin{verbatim}
calculate the total of all the items in the list.

$item = array( 'total' => $item, 'total' => $item, 'total'
=> $item, '
\end{verbatim}
Con \texttt{do\_sample=True} y \texttt{top\_k=1}, las 25 corridas registradas en el CSV (5 ejecuciones del notebook con 5 semillas cada una) produjeron 1 salida distinta: idéntica a la de greedy (Tabla~\ref{tab:p0k1}).\begin{table}[H]\centering\caption{Parte 0.b: cinco corridas con \texttt{do\_sample=True} y \texttt{top\_k=1}. IDs completos en \texttt{datos/parte0\_generaciones.jsonl}; soporte/paso = tokens con probabilidad no nula tras el filtro.}\label{tab:p0k1}
{\footnotesize\setlength{\tabcolsep}{3.5pt}
\begin{tabular}{@{}llp{4.6cm}lll@{}}
\hline
corrida & semilla & primeros 10 IDs & tokens & soporte/paso & = greedy \\
\hline
0 & 123 & 15284 262 2472 286 477 262 3709 287 262 1351 ... & 40 & 1 & sí \\
1 & 124 & 15284 262 2472 286 477 262 3709 287 262 1351 ... & 40 & 1 & sí \\
2 & 125 & 15284 262 2472 286 477 262 3709 287 262 1351 ... & 40 & 1 & sí \\
3 & 126 & 15284 262 2472 286 477 262 3709 287 262 1351 ... & 40 & 1 & sí \\
4 & 127 & 15284 262 2472 286 477 262 3709 287 262 1351 ... & 40 & 1 & sí \\
\hline
\end{tabular}
}\end{table}


Degeneración con 100 tokens en modo greedy (Holtzman et al.~\cite{holtzman2020}), salida sin editar:
\begin{verbatim}
calculate the total of all the items in the list.

$item = array( 'total' => $item, 'total' => $item, 'total'
=> $item, 'total' => $item, 'total' => $item, 'total' =>
$item, 'total' => $item, 'total' => $item, 'total' =>
$item, 'total' => $item, 'total' => $item, 'total' => $
\end{verbatim}

\begin{figure}[H]\centering\includegraphics[width=0.95\textwidth,height=0.42\textheight,keepaspectratio]{figuras/parte0b_cortes.png}\caption{Dónde corta cada palanca: tokens que sobreviven a top\_k=5 y a top\_p=0.9 en cada prefijo. En el prefijo abierto los dos conjuntos no coinciden (5 frente a 241 tokens).}\label{fig:p0b}\end{figure}


\subsection{El límite del modelo base}
Salida cruda de GPT-2 ante la instrucción exacta de las Partes 1 a 4:
\begin{verbatim}
Una cada una de 43 sillas. ¿Cuántos tornillos le sobran?

Una cada una de 43 sillas. ¿Cuántos tornillos le sobran?

Una cada una de 43 sillas. ¿Cuántos tornillos le sobran?

Una cada una de 43 sillas. ¿Cuántos
\end{verbatim}
GPT-2 no entrega el entero pedido: copia un fragmento del enunciado («Una cada
una de 43 sillas. ¿Cuántos tornillos le sobran?») y lo repite hasta agotar los
100 tokens, el mismo bucle de la degeneración de 0.b. No intenta calcular
porque para él no existe la noción de pregunta: fue entrenado solo para
predecir el siguiente token de un corpus, y lo que sigue estadísticamente a un
problema escrito es más texto parecido. Se suma que el prompt está en español y su corpus es casi todo inglés. Lo que tienen los modelos de
las Partes 1 a 4 y a este le falta es la etapa de alineación posterior al
preentrenamiento, que enseña a tratar la entrada como una orden y a producir
la forma de respuesta esperada. Esa etapa no garantiza que el número sea correcto; el verificador sigue siendo necesario.

\section{Parte 1 --- Comparación de modelos}
Precios de la tabla semestral, cada uno con la fecha de verificación de su propia fila (Tabla~\ref{tab:p1}).
\begin{table}[H]\centering\caption{Parte 1: tres modelos sobre los 10 casos. Latencia = tiempo de pared de la petición.}\label{tab:p1}
{\footnotesize\setlength{\tabcolsep}{3.5pt}
\begin{tabular}{lrrrrrl}
\hline
modelo & exactitud & lat. (s) & tok. in & tok. out & costo USD & verificado \\
\hline
gpt-4o-mini & 0 & 1.9063 & 793 & 12 & 0.00013 & 2026-08-26 \\
gpt-5.6-luna & 1 & 1.13 & 783 & 405 & 0.00064 & 2026-09-18 \\
qwen3:1.7b & 0.9 & 10.677 & 1011 & 7424 & 0 & local \\
\hline
\end{tabular}
}\end{table}
Elegiríamos gpt-5.6-luna: acertó los diez casos, fue el más rápido (1,13 s
frente a 1,91 s de gpt-4o-mini y 10,7 s de qwen3 en local) y costó 0,00064 USD
por los diez. Qwen3 acertó nueve a costo cero, pero necesita GPU y tarda porque
razona antes de responder; lo consideraríamos para trabajo local o con datos
que no puedan salir de la máquina. Observaciones cualitativas: gpt-4o-mini
respetó el formato (un solo número, sin unidades) en los diez casos, no fue
verboso (1,2 tokens de salida por llamada) y falló los diez: dio cifras
equivocadas con seguridad, sin señal de duda. gpt-5.6-luna respetó el formato
y no mostró errores ni respuestas vacías. qwen3 respetó el formato en la
respuesta visible, pero generó 7 424 tokens de salida en diez casos, casi todos
de razonamiento; su único fallo fue de cálculo, no de formato. Ninguno devolvió
errores de API. La latencia es tiempo de pared de la petición completa, red
incluida. Ampliaríamos el conjunto antes de decidir en producción: diez casos
permiten comparar, no generalizar.

Confirmamos que los precios se leyeron de la tabla semestral del enunciado,
con la fecha de cada fila, y no de memoria.

\section{Parte 2 --- Decodificación}
\subsection{Matriz de exposición}
Cada parámetro se mandó con dos valores extremos, 3 casos $\times$ 2 corridas. El estado se decide por la dispersión entre corridas de la tupla (salida, tokens de salida), no por la respuesta final (Tabla~\ref{tab:p2a}); los mensajes de error se transcriben completos a continuación de la tabla.
\begin{table}[H]\centering\caption{Parte 2.a: matriz de exposición, declarado frente a observado (2026-09-20). Evidencia: porcentaje de casos cuyas corridas difieren con cada valor; o número de llamadas con error.}\label{tab:p2a}
{\footnotesize\setlength{\tabcolsep}{3.5pt}
\begin{tabular}{@{}lllp{2.0cm}p{3.3cm}@{}}
\hline
modelo & parámetro & declarado (tabla) & observado & evidencia \\
\hline
gpt-4o-mini & temperature & si & aceptado y actua & temperature=0: 0\% / temperature=1.5: 100\% \\
gpt-4o-mini & top\_p & si & aceptado y actua & top\_p=0.1: 33\% / top\_p=1: 67\% \\
gpt-4o-mini & top\_k & no & rechazado & 12/12 llamadas con error \\
gpt-5.6-luna & temperature & no\_en\_esta\_fila & rechazado & 12/12 llamadas con error \\
gpt-5.6-luna & top\_p & no\_en\_esta\_fila & rechazado & 6/12 llamadas con error \\
gpt-5.6-luna & top\_k & no & rechazado & 12/12 llamadas con error \\
qwen3:1.7b & temperature & si & aceptado y actua & temperature=0: 33\% / temperature=1.5: 100\% \\
qwen3:1.7b & top\_p & si & aceptado y actua & top\_p=0.1: 33\% / top\_p=1: 100\% \\
qwen3:1.7b & top\_k & si & aceptado y actua & top\_k=1: 33\% / top\_k=40: 100\% \\
\hline
\end{tabular}
}\end{table}
Mensajes literales completos de las celdas rechazadas:
\smallskip\noindent\texttt{gpt-4o-mini} / \texttt{top\_k}:
\begin{verbatim}
TypeError: Completions.create() got an unexpected keyword argument
'top_k'
\end{verbatim}
\smallskip\noindent\texttt{gpt-5.6-luna} / \texttt{temperature}:
\begin{verbatim}
400: Error code: 400 - {'error': {'message': "Unsupported value:
'temperature' does not support 0.0 with this model. Only the
default (1) value is supported.", 'type': 'invalid_request_error',
'param': 'temperature', 'code': 'unsupported_value'}}
\end{verbatim}
\smallskip\noindent\texttt{gpt-5.6-luna} / \texttt{top\_p}:
\begin{verbatim}
400: Error code: 400 - {'error': {'message': "Unsupported
parameter: 'top_p' is not supported with this model.", 'type':
'invalid_request_error', 'param': 'top_p', 'code':
'unsupported_parameter'}}
\end{verbatim}
\smallskip\noindent\texttt{gpt-5.6-luna} / \texttt{top\_k}:
\begin{verbatim}
TypeError: Completions.create() got an unexpected keyword argument
'top_k'
\end{verbatim}
En gpt-4o-mini, temperatura y top-p mostraron cambios entre corridas. Top-k
falló en el SDK, antes de llegar al servidor. Luna rechazó las temperaturas
probadas y aceptó top\_p=1, pero rechazó 0.1. En Ollama observamos contenido en
thinking sin pedirlo, contrario al valor predeterminado indicado en la tabla.

\subsection{Barrido de temperature y top-p}
Estimación previa con el supuesto del enunciado (1\,000 tokens de entrada y 600 de salida por llamada) al precio de gpt-4o-mini verificado el 2026-08-26: 750 llamadas $\approx$ 0.38\,USD. Se corrieron las 750 (15 celdas $\times$ 5 corridas $\times$ 10 casos); costo real 0.033\,USD, porque la salida media fue de 1--2 tokens y no de 600.
\begin{table}[H]\centering\caption{Parte 2.b: rejilla temperature $\times$ top\_p sobre gpt-4o-mini, 5 corridas $\times$ 10 casos por celda (750 llamadas).}\label{tab:p2b}
{\footnotesize\setlength{\tabcolsep}{3.5pt}
\begin{tabular}{rrrrr}
\hline
T & top\_p & exactitud & tokens out (media) & estabilidad \\
\hline
0.000 & 0.500 & 0.000 & 1.200 & 0.960 \\
0.000 & 0.900 & 0.000 & 1.200 & 0.960 \\
0.000 & 1.000 & 0.000 & 1.240 & 0.980 \\
0.300 & 0.500 & 0.000 & 1.340 & 0.740 \\
0.300 & 0.900 & 0.020 & 1.240 & 0.820 \\
0.300 & 1.000 & 0.000 & 1.240 & 0.740 \\
0.700 & 0.500 & 0.000 & 1.360 & 0.640 \\
0.700 & 0.900 & 0.020 & 1.180 & 0.560 \\
0.700 & 1.000 & 0.000 & 1.240 & 0.560 \\
1.000 & 0.500 & 0.000 & 1.280 & 0.640 \\
1.000 & 0.900 & 0.000 & 1.220 & 0.420 \\
1.000 & 1.000 & 0.000 & 1.340 & 0.480 \\
1.500 & 0.500 & 0.000 & 1.260 & 0.580 \\
1.500 & 0.900 & 0.020 & 1.340 & 0.440 \\
1.500 & 1.000 & 0.020 & 771.140 & 0.280 \\
\hline
\end{tabular}
}\end{table}
\begin{figure}[H]\centering\includegraphics[width=0.6\textwidth]{figuras/parte2b_rejilla.png}\caption{Exactitud por celda de la rejilla (gpt-4o-mini).}\label{fig:p2b}\end{figure}
\begin{table}[H]\centering\caption{Parte 2.b: barrido de top\_k en local (qwen3:1.7b, T=1, 5 corridas secuenciales $\times$ 10 casos).}\label{tab:p2bk}
{\footnotesize\setlength{\tabcolsep}{3.5pt}
\begin{tabular}{rrrrr}
\hline
top\_k & exactitud & tokens out (media) & estabilidad & casos idénticos x5 \\
\hline
1.000 & 0.920 & 1080.320 & 0.960 & 8 \\
5.000 & 0.960 & 947.780 & 0.880 & 7 \\
40.000 & 0.980 & 931.120 & 0.860 & 7 \\
\hline
\end{tabular}
}\end{table}
La temperatura divide los logits antes del softmax: con T<1 amplía la
distancia entre el token más probable y el resto y concentra la masa; con T>1
la reduce y reparte masa hacia la cola. top\_p corta esa cola por masa
acumulada. Sobre gpt-4o-mini (750 llamadas), la exactitud fue 0 \% en once celdas y 2 \% (una de cincuenta) en las
otras cuatro: el muestreo no cambia lo que el modelo sabe hacer, solo cuánto
varía su salida. Lo que sí se mueve es la estabilidad: entre 0,96 y 0,98 con
T=0, alrededor de 0,75 con T=0,3 y 0,28 en T=1,5 con top\_p=1. Esa celda además
degenera: 771 tokens de salida en promedio frente a 1,2-1,4 en el resto,
sin corte de núcleo la cola aplanada no termina; con
top\_p=0,5 la misma T=1,5 vuelve a 1,26 tokens. Que T=0 no dé 1,0 es un hallazgo: la rama argmax no es una promesa de
reproducibilidad en una API. Para esta tarea elegiríamos T=0 y
top\_p=1, por estabilidad, ya que ninguna celda mejora la exactitud: la única palanca que la mueve es el prompt (Parte 3) o el modelo.
Para generar tres titulares alternativos elegiríamos T=1 con top\_p=0,9: hay
dispersión entre corridas y el núcleo excluye continuaciones absurdas. En local, top\_k=1 dio la misma salida en
ocho de diez casos en cinco corridas secuenciales, no en los diez, y la
dispersión creció con k (estabilidad 0,96 → 0,86 de top\_k=1 a 40) sin cambiar
la exactitud (92-98 \%).

\section{Parte 3 --- Prompting estructurado}
Las cuatro plantillas están en \texttt{src/prompts.py}. Los ejemplos del few-shot son problemas inventados, no los diez casos. La variante estructurada usa el modo JSON del proveedor (\texttt{response\_format} con tipo \texttt{json\_object}), sin esquema estricto.
\begin{table}[H]\centering\caption{Parte 3: cuatro variantes de prompting sobre gpt-4o-mini, 10 casos.}\label{tab:p3}
{\footnotesize\setlength{\tabcolsep}{3.5pt}
\begin{tabular}{lrrrr}
\hline
variante & exactitud & tokens in & tokens out & costo USD \\
\hline
zero\_shot & 0.1 & 793 & 11 & 0.00013 \\
few\_shot & 0 & 1842 & 11 & 0.00028 \\
cot & 1 & 933 & 3160 & 0.00204 \\
structured & 0 & 883 & 61 & 0.00017 \\
\hline
\end{tabular}
}\end{table}

Respuestas completas de cada variante sobre el mismo caso (las diez por variante están en el CSV, columna \texttt{salida}):


\smallskip\noindent\textbf{zero\_shot} (caso c01, esperado 311, extraído 245.0, acierto no):
\begin{verbatim}
245
\end{verbatim}


\smallskip\noindent\textbf{few\_shot} (caso c01, esperado 311, extraído 81.0, acierto no):
\begin{verbatim}
81
\end{verbatim}


\smallskip\noindent\textbf{cot} (caso c01, esperado 311, extraído 311.0, acierto sí):
\begin{verbatim}
Para resolver este problema, primero vamos a determinar
cuántos tornillos hay en las cajas y luego cuántos
tornillos se usan para fabricar las sillas.

1. **Calcular el total de tornillos comprados:**
   - El taller compra 17 cajas.
   - Cada caja contiene 36 tornillos.
   - Por lo tanto, el total de tornillos comprados es:
     \[
     17 \text{ cajas} \times 36 \text{ tornillos/caja} =
612 \text{ tornillos}
     \]

2. **Calcular cuántos tornillos se usan para las sillas:**
   - Se usan 7 tornillos para cada silla.
   - El taller necesita hacer 43 sillas.
   - Por lo tanto, el total de tornillos usados es:
     \[
     7 \text{ tornillos/silla} \times 43 \text{ sillas} =
301 \text{ tornillos}
     \]

3. **Calcular cuántos tornillos sobran:**
   - Para encontrar cuántos tornillos sobran, restamos los
tornillos usados del total de tornillos comprados:
     \[
     612 \text{ tornillos} - 301 \text{ tornillos} = 311
\text{ tornillos}
     \]

Finalmente, el número de tornillos que le sobran al taller
es:

Respuesta: 311
\end{verbatim}


\smallskip\noindent\textbf{structured} (caso c01, esperado 311, extraído 71.0, acierto no):
\begin{verbatim}
{"respuesta": 71}
\end{verbatim}
Zero-shot acertó un caso y few-shot ninguno. Las dos piden un número directo y
gpt-4o-mini no puede encadenar dos o tres operaciones en un solo token; los
ejemplos duplican los tokens de entrada (1 842 frente a 793 en los diez casos)
sin ganar nada, porque enseñan el formato, no el procedimiento.
Chain-of-Thought pasó a diez de diez a cambio de 3 160 tokens de salida, casi
300 veces los 11 de zero-shot, y un costo 16 veces mayor (0,00204 frente a
0,00013 USD): compra cómputo intermedio, que es lo que esta tarea necesita. La
variante estructurada usó el modo JSON del proveedor (response\_format con
json\_object), sin esquema estricto: las diez salidas parsearon como JSON y
ninguna fue correcta. Distinguimos lo que garantiza el proveedor de lo que
verificamos nosotros: el modo JSON garantiza la forma; el valor lo comprobó
nuestro verificador. Usaríamos CoT cuando la tarea tenga pasos dependientes y
la exactitud pese más que el costo y la latencia; zero-shot cuando la
respuesta sea de recuperación directa o una etiqueta, donde razonar no añade
información; few-shot cuando lo difícil sea el formato o el criterio de
clasificación, no el cálculo; y salida estructurada siempre que otro programa
consuma la respuesta, combinada con CoT si el contenido lo exige, porque por
sí sola solo garantiza que parsea. No asumiríamos que combinar técnicas mejora
el resultado sin medir esa combinación con los mismos diez casos.

\section{Parte 4 --- Modelos de razonamiento y niveles de esfuerzo}
\subsection{Barrido de esfuerzo}
Modelo gpt-5.6-luna. El dial es \texttt{reasoning\_effort}. El contador de tokens de razonamiento es el campo \texttt{reasoning\_tokens} que devuelve la API. Se corrió \texttt{low} completo antes de los otros dos: 388 tokens de salida totales en 10 casos (38,8 por llamada, razonamiento incluido) frente a los 1\,500 por llamada que supone el presupuesto del enunciado para el nivel bajo; la cuenta se rehizo con esa cifra antes de seguir.
\begin{table}[H]\centering\caption{Parte 4.a: barrido de esfuerzo en gpt-5.6-luna, 10 casos por nivel.}\label{tab:p4a}
{\footnotesize\setlength{\tabcolsep}{3.5pt}
\begin{tabular}{lrrrrr}
\hline
esfuerzo & exactitud & tok. razonamiento & tok. visibles & lat. (s) & costo USD \\
\hline
low & 1 & 287 & 101 & 1.2317 & 0.00062 \\
medium & 1 & 288 & 101 & 0.95785 & 0.00062 \\
high & 1 & 375 & 101 & 1.0125 & 0.00073 \\
\hline
\end{tabular}
}\end{table}
\begin{figure}[H]\centering\includegraphics[width=0.62\textwidth]{figuras/parte4a_exactitud_vs_tokens.png}\caption{Exactitud frente a tokens de razonamiento medidos.}\label{fig:p4a1}\end{figure}
\begin{figure}[H]\centering\includegraphics[width=0.62\textwidth]{figuras/parte4a_costo_vs_exactitud.png}\caption{Costo en USD frente a exactitud.}\label{fig:p4a2}\end{figure}
La curva es plana: diez de diez en low, medium y high. Los tokens de
razonamiento medidos fueron 287, 288 y 375 en los diez casos. Low y medium son
indistinguibles en el contador, así que entre esos dos niveles el dial se
comporta como «aceptado y no actúa», el mismo estado del medio de la Parte
2.a, diagnosticado con el contador y no con la prosa de la respuesta. High sí
mueve el contador (+30 \%) y el costo, de 0,00062 a 0,00073 USD, sin ganar un
punto de exactitud: pagamos 0,0001 USD adicionales por cero puntos, así que no
podemos calcular un costo por punto. El punto de rendimientos decrecientes
está por debajo de low para este conjunto, y no podemos ubicarlo con precisión
porque no evaluamos niveles inferiores. Las cifras también corrigen el
supuesto del presupuesto: 38,8 tokens de salida totales por llamada en low
frente a los 1 500 que supone el enunciado, por lo que el costo real fue unas
40 veces menor. La lección no es que el dial no sirva, sino que el nivel se
elige por tarea y se mide: estos diez casos están por debajo del umbral de
dificultad donde el esfuerzo empieza a pagar. Con casos más difíciles y varias
corridas por nivel esperaríamos ver la curva.

\subsection{El caso donde pensar más hace daño}
Casos c11 (distracción por lo irrelevante), c12 (sobreajuste al marco) y c13 (correlación espuria), según los modos de fallo de Gema et al.~\cite{gema2025}.
\begin{table}[H]\centering\caption{Parte 4.b: los tres casos contaminados en los niveles extremos.}\label{tab:p4b}
{\footnotesize\setlength{\tabcolsep}{3.5pt}
\begin{tabular}{llrr}
\hline
esfuerzo & caso & acierto & tokens razonamiento \\
\hline
high & c11 & 1 & 0 \\
high & c12 & 1 & 24 \\
high & c13 & 1 & 26 \\
low & c11 & 1 & 0 \\
low & c12 & 1 & 0 \\
low & c13 & 1 & 21 \\
\hline
\end{tabular}
}\end{table}
No observamos empeoramiento: ambos niveles acertaron los tres casos. No reprodujimos
los fallos descritos por Gema et al., pero eso no contradice su estudio. Tres
casos no permiten establecer una regla general.

\section{Parte 5 --- Reflexión teórica}
\subsection{Explica intuitivamente qué hace el mecanismo de atención en un transformer. ¿Por qué es importante el concepto de query, key, value?}
Entendemos la atención como una forma de decidir qué partes del contexto
aportan información para representar cada token. Cada posición produce tres
vectores por proyección lineal: la query expresa lo que esa posición busca, la
key permite que las demás posiciones se comparen con esa búsqueda y el value
contiene la información que se mezcla. El modelo calcula puntuaciones entre la
query y todas las keys, las convierte en pesos con softmax y obtiene una
combinación ponderada de los values. Así relaciona palabras cercanas o
distantes según el contexto, y todas las posiciones se calculan en paralelo,
sin el cuello de botella secuencial de una red recurrente. Separar los tres
roles permite aprender por separado qué comparar y qué contenido recuperar. En
generación causal, una máscara impide atender a tokens futuros.

\subsection{¿Cuál es la diferencia entre un modelo base y un modelo alineado? ¿Por qué un modelo base no es directamente útil para las aplicaciones de este curso?}
Salida completa de 0.c ante la instrucción de la tarea, sin editar:\begin{verbatim}
Una cada una de 43 sillas. ¿Cuántos tornillos le sobran?

Una cada una de 43 sillas. ¿Cuántos tornillos le sobran?

Una cada una de 43 sillas. ¿Cuántos tornillos le sobran?

Una cada una de 43 sillas. ¿Cuántos
\end{verbatim}En nuestra prueba, GPT-2 continuó el enunciado sin resolverlo. Lo relacionamos con la diferencia entre
continuar texto y estar ajustado para seguir instrucciones. Un modelo base
aprende patrones del corpus durante el preentrenamiento y solo predice el
siguiente token. Un modelo alineado recibe después ejemplos de peticiones y
respuestas, y a menudo entrenamiento con preferencias humanas; eso le enseña a
tratar la entrada como una orden y a responder en la forma esperada. Para las
aplicaciones del curso necesitamos que la salida responda a una tarea y
respete un formato, algo que no podemos dar por hecho con el modelo base. Aun
así, el modelo base nos resultó útil para observar directamente los logits y
el muestreo. El ajuste mejora el comportamiento, pero no elimina la necesidad
de verificar las respuestas.

\subsection{En la Parte 3, la variante con Chain-of-Thought probablemente usó más tokens que la zero-shot. ¿Cuándo vale la pena pagar ese costo y cuándo no?}
En nuestros resultados, CoT pasó de un acierto a diez frente a zero-shot, con
los tokens de salida de 11 a 3 160 y el costo de 0,00013 a 0,00204 USD por los
diez casos. Pagaríamos esa diferencia cuando la respuesta dependa de pasos
intermedios que el modelo no resuelve en un solo token, cuando el costo de un
error supere con mucho el de los tokens y cuando la latencia adicional sea
aceptable. No la pagaríamos para extraer un dato explícito o devolver una
etiqueta, donde primero probaríamos una instrucción corta; ni en aplicaciones
de alto volumen con margen pequeño por llamada. También compararíamos con un
modelo de razonamiento: en la Parte 4, gpt-5.6-luna resolvió los mismos casos
con 38,8 tokens de salida por llamada en low, frente a 316 de CoT en
gpt-4o-mini. Gastar más tokens solo tiene sentido si la mejora medida lo
justifica.

\subsection{En tu matriz de la Parte 2.a hay al menos una celda «rechazado». ¿Qué gana un proveedor al retirar un parámetro de muestreo de su API, y qué pierde quien lo usaba?}
En nuestras pruebas, gpt-5.6-luna rechazó con 400 los valores de temperature y
top\_p distintos del predeterminado. Consideramos que el proveedor gana control
sobre la distribución de salida de un modelo cuyo comportamiento depende de
una fase de razonamiento: con muestreo libre el usuario podría degradar esa
fase y el proveedor tendría que dar soporte a resultados que no puede
reproducir. Es una explicación plausible, no una intención
que podamos demostrar con el CSV. Quien usaba la palanca pierde la capacidad
de intercambiar exactitud por diversidad en la misma llamada y la
reproducibilidad aproximada de T=0. Lo que la sustituye es el nivel de
esfuerzo: en vez de controlar cómo se muestrea, se controla cuánto se piensa.
La diversidad, cuando hace falta, se obtiene pidiendo varias alternativas en
el prompt o con varias llamadas, y no equivale a la temperatura.

\subsection{Los tokens de razonamiento de la Parte 4 se facturan y casi nunca se devuelven en crudo. Si tu modelo te mostrara la traza completa, ¿sería eso una explicación de cómo llegó a la respuesta?}
No tomaríamos una traza completa como explicación suficiente de cómo se obtuvo
la respuesta. Lanham et al. (2023) miden si la respuesta cambia cuando se
altera la cadena de razonamiento y encuentran que en muchos casos el modelo
llega al mismo resultado con una traza truncada o con errores insertados, y
que esa fidelidad varía con el tamaño del modelo y con la tarea. Una
explicación puede ser coherente sin reflejar el proceso que produjo la
respuesta. Para auditar un sistema en producción, esto significa que la traza
es evidencia de proceso, no prueba de causa: sirve para detectar errores
visibles y para medir cuánto cómputo se gastó, como el contador de la Parte 4,
pero no basta para certificar por qué el modelo respondió lo que respondió. La
auditoría debe apoyarse en pruebas de comportamiento sobre casos verificables y registro de versiones y condiciones de
ejecución.



\section{Reproducibilidad y desviaciones declaradas}
\begin{itemize}
\item Código, casos, \texttt{requirements.txt} y \texttt{datos/resultados.csv} (1159 filas, una por llamada, las locales con costo cero) en \url{https://github.com/jastk45/Maestria-}.\newline Carpeta: \path{IA-Generativa-y-Agentes/semana1/taller01_foundation_models}
\item Las claves se leen de un \texttt{.env} excluido por \texttt{.gitignore}; ninguna aparece en el informe, el notebook ni el repositorio.
\item Las sondas de 2.a y la rejilla de 2.b se mandaron con 4 peticiones en paralelo (por eso $T=0$ no da estabilidad 1,0); el barrido de top\_k en local, en secuencia.
\item Los tokens de razonamiento de qwen3 se miden en tokens: \texttt{eval\_count} de Ollama (todo lo generado) menos los tokens de la respuesta visible contados con el tokenizador de Qwen3. Los de gpt-5.6-luna salen del contador \texttt{reasoning\_tokens} de la API.
\end{itemize}

\begin{thebibliography}{8}
\bibitem{holtzman2020} Holtzman, A., Buys, J., Du, L., Forbes, M., Choi, Y.: The Curious Case of Neural Text Degeneration. ICLR (2020). \url{https://arxiv.org/abs/1904.09751}
\bibitem{gema2025} Gema, A.P., et al.: Inverse Scaling in Test-Time Compute. TMLR (2025)
\bibitem{lanham2023} Lanham, T., et al.: Measuring Faithfulness in Chain-of-Thought Reasoning. arXiv:2307.13702 (2023)
\bibitem{vaswani2017} Vaswani, A., et al.: Attention Is All You Need. NeurIPS (2017)
\bibitem{gpt2} Radford, A., et al.: Language Models are Unsupervised Multitask Learners. OpenAI (2019). Model card: \url{https://huggingface.co/openai-community/gpt2}
\end{thebibliography}
\end{document}
